
\documentclass{article}
\usepackage{colortbl}
\usepackage{makecell}
\usepackage{multirow}
\usepackage{supertabular}

\begin{document}

\newcounter{utterance}

\centering \large Interaction Transcript for game `codenames', experiment `opponent\_none', episode 1 with Apertus{-}v1.5{-}8B.
\vspace{24pt}

{ \footnotesize  \setcounter{utterance}{1}
\setlength{\tabcolsep}{0pt}
\begin{supertabular}{c@{$\;$}|p{.15\linewidth}@{}p{.15\linewidth}p{.15\linewidth}p{.15\linewidth}p{.15\linewidth}p{.15\linewidth}}
    \# & $\;$A & \multicolumn{4}{c}{Game Master} & $\;\:$B\\
    \hline

    \theutterance \stepcounter{utterance}  
    & & & \multicolumn{2}{p{0.3\linewidth}}{
        \cellcolor[rgb]{0.9,0.9,0.9}{
            \makecell[{{p{\linewidth}}}]{
                \texttt{\tiny{[GM$|$GM]}}
                {'hidden': {'team': ['jama', 'antarktída', 'bomba', 'citrón', 'peking', 'oliva', 'pilot', 'prístav', 'mikroskop'], 'innocent': ['trať', 'čiapka', 'model', 'zmena', 'karta', 'palica', 'hviezda'], 'opponent': ['ruka', 'sila', 'fénix', 'kravata', 'hodinky', 'kráľ', 'tekvica', 'dlhopis'], 'assassin': ['papier']}, 'revealed': {'team': {'team': [], 'innocent': [], 'opponent': [], 'assassin': []}, 'opponent': {'team': [], 'innocent': [], 'opponent': [], 'assassin': []}}}
            }
        }
    }
    & & \\ \\

    \theutterance \stepcounter{utterance}  
    & & \multicolumn{4}{p{0.6\linewidth}}{
        \cellcolor[rgb]{0.9,0.9,0.9}{
            \makecell[{{p{\linewidth}}}]{
                \texttt{\tiny{[P1$\langle$GM]}}
                \texttt{Uveďte jednu jednoslovnú nápovedu, ktorá súvisí s niektorými slovami vášho tímu z nižšie uvedeného zoznamu. Môžete si vybrať, na koľko a ktoré slová sa vaša nápoveda vzťahuje.} \\
\texttt{Cieľom je, aby váš partner uhádol všetky slová vášho tímu na základe nápovied, ktoré mu poskytnete, skôr ako druhý tím odhalí všetky svoje. Ak váš partner uhádne slovo "súpera" alebo "rozptyľujúce" slovo, jeho kolo hádania sa končí a na rad prichádza súperov tím. Ak váš partner uhádne slovo "zabiť", hra sa končí a váš tím okamžite prehráva.} \\
\texttt{Čím viac slov sa týka vaša nápoveda, tým rýchlejšie môžete vyhrať hru, ale tým ťažšie môžu byť asociácie na uhádnutie pre vášho partnera, preto si nápovedné slovo vyberte múdro. Nápovedné slovo musí byť sémanticky spojené s cieľovými slovami, nemôže byť jedným zo slov v zoznamoch ani obsahovať ich časti.} \\
\\ 
\texttt{Svoju jednoslovnú nápovedu a zoznam príbuzných cieľových slov oddelených čiarkou uvádzajte vždy v nasledujúcom formáte a odpovede robte čo najkratšie, nikdy neuvádzajte žiadny iný text, ako je požadovaný v tomto formáte:} \\
\\ 
\texttt{PORADIE: <SLOVO>} \\
\texttt{CIELE: <CIELE>} \\
\\ 
\texttt{Slová vášho tímu sú: jama, antarktída, bomba, citrón, peking, oliva, pilot, prístav, mikroskop.} \\
\texttt{Slová súperovho tímu sú: ruka, sila, fénix, kravata, hodinky, kráľ, tekvica, dlhopis.} \\
\texttt{Rozptyľujúce slová sú: trať, čiapka, model, zmena, karta, palica, hviezda.} \\
\texttt{Zabíjacie slová sú: papier.} \\
            }
        }
    }
    & \\ \\

    \theutterance \stepcounter{utterance}  
    & \multicolumn{4}{p{0.6\linewidth}}{
        \cellcolor[rgb]{0.9,0.9,0.9}{
            \makecell[{{p{\linewidth}}}]{
                \texttt{\tiny{[P1$\rangle$GM]}}
                \texttt{PORADIE: voda} \\
\texttt{CIELE: citrón, oliva} \\
            }
        }
    }
    & & \\ \\

    \theutterance \stepcounter{utterance}  
    & & & \multicolumn{2}{p{0.3\linewidth}}{
        \cellcolor[rgb]{0.9,0.9,0.9}{
            \makecell[{{p{\linewidth}}}]{
                \texttt{\tiny{[GM$|$GM]}}
                \texttt{voda} \\
            }
        }
    }
    & & \\ \\

    \theutterance \stepcounter{utterance}  
    & & & \multicolumn{2}{p{0.3\linewidth}}{
        \cellcolor[rgb]{0.9,0.9,0.9}{
            \makecell[{{p{\linewidth}}}]{
                \texttt{\tiny{[GM$|$GM]}}
                ['citrón', 'oliva']
            }
        }
    }
    & & \\ \\

    \theutterance \stepcounter{utterance}  
    & & & \multicolumn{2}{p{0.3\linewidth}}{
        \cellcolor[rgb]{0.9,0.9,0.9}{
            \makecell[{{p{\linewidth}}}]{
                \texttt{\tiny{[GM$|$GM]}}
                {'word': 'citrón', 'assignment': 'team'}
            }
        }
    }
    & & \\ \\

    \theutterance \stepcounter{utterance}  
    & & & \multicolumn{2}{p{0.3\linewidth}}{
        \cellcolor[rgb]{0.9,0.9,0.9}{
            \makecell[{{p{\linewidth}}}]{
                \texttt{\tiny{[GM$|$GM]}}
                {'word': 'oliva', 'assignment': 'team'}
            }
        }
    }
    & & \\ \\

    \theutterance \stepcounter{utterance}  
    & & \multicolumn{4}{p{0.6\linewidth}}{
        \cellcolor[rgb]{0.9,0.9,0.9}{
            \makecell[{{p{\linewidth}}}]{
                \texttt{\tiny{[GM$\rangle$P2]}}
                \texttt{Uveďte zoznam maximálne 2 slov z nasledujúceho zoznamu oddelených čiarkou, ktoré najlepšie súvisia so slovom 'voda' alebo sú s ním najviac prepojené. Zoznam hádaní vždy začnite slovom HÁDANIE: a do odpovede neuvádzajte žiadny iný text.} \\
\\ 
\texttt{sila, antarktída, ruka, palica, prístav, jama, oliva, bomba, zmena, fénix, dlhopis, model, kráľ, čiapka, citrón, hodinky, kravata, peking, trať, papier, pilot, karta, hviezda, mikroskop, tekvica} \\
            }
        }
    }
    & \\ \\

    \theutterance \stepcounter{utterance}  
    & & & \multicolumn{4}{p{0.6\linewidth}}{
        \cellcolor[rgb]{0.9,0.9,0.9}{
            \makecell[{{p{\linewidth}}}]{
                \texttt{\tiny{[GM$\langle$P2]}}
                \texttt{HÁDANIE: voda, prístav} \\
            }
        }
    }
     \\ \\

    \theutterance \stepcounter{utterance}  
    & & & \multicolumn{2}{p{0.3\linewidth}}{
        \cellcolor[rgb]{0.9,0.9,0.9}{
            \makecell[{{p{\linewidth}}}]{
                \texttt{\tiny{[GM$|$GM]}}
                {'player': 'Guesser', 'type': 'guess is clue word', 'utterance': 'voda, prístav', 'message': "Guessed word 'voda' is the same word as the provided clue word, you should only select words from the provided list.", 'guess': 'voda', 'clue': 'voda'}
            }
        }
    }
    & & \\ \\

\end{supertabular}
}

\end{document}
